The toolkit contains a couple of functions for generating the grids and transition matrices for exogenous variables based on the Tauchen Method for
discrete state space approximations of VARs. The main command $TauchenMethod$ is for creating the approximation directly from the parameters of an AR(1),
and this is now described.

The Toolkit also includes commands $TauchenMethodVAR$ for approximation (quadrature) of a VAR(1), and $RouwenhorstMethod$ which provides an alternative 
approximations of AR(1) processes that outperforms the Tauchen method when the autocorrelation coefficient (referred to below as $\rho$) takes a value
near one.


\noindent Namely,
\begin{itemize}
%  \item \textit{TauchenMethod\_{}Data}: Generates a markov chain approximation from the time series input (uses Tauchen method to approximate an AR(1) process
%       estimated from that time series).
 \item \textit{TauchenMethod}: Generates a markov chain approximation for AR(1) process defined by the inputed parameters.
\end{itemize}

% Since the quality of the approximation depends on 'q', one thing you can do is use \textit{MarkovChainMoments (undocumented)} to check how closely the Markov process generated
% by the Tauchen method approximates the parameters you used to generate it (note: should you so wish you could even search across q's to find the one that 
% gives the best approximation, in terms of matching the three parameters, given the number of states you are using).

\subsection{TauchenMethod}
To create a discrete state space approximation of the AR(1) process,
\begin{equation} \label{eqn:ar1}
 z_t =\mu +\rho z_{t-1} +\epsilon_t \quad \epsilon_t \sim^{iid} N(0,\sigma^2)
\end{equation}
the command is, \\
\indent \textcolor{RoyalBlue}{[z\_{}grid, pi\_{}z]=TauchenMethod(mu,sigmasq,rho,znum,q, [tauchenoptions]);}
the outputs are the grid of discretized values ($z\_{}grid$) and the transition matrix ($pi\_{}z$). 
We now describe in more detail the inputs and outputs.

\subsubsection{Inputs and Outputs}
The inputs are
\begin{itemize}
 \item Define the constant, autocorrelation, and error variance terms relating to equation \ref{eqn:ar1}, eg. \\
    \textcolor{RoyalBlue}{mu=1; rho=0.9; sigmasq=0.09;}
 \item Set the number of points to use for the approximation, \\
    \textcolor{RoyalBlue}{znum = 9;} \\
    (When using the value function iteration commands, $znum$ will correspond to $n\_{}z$ whenever it is the only shock.)
 \item Set $q$, roughly q defines the number of standard deviations of $z$ covered by the range of discrete points, \\
   \textcolor{RoyalBlue}{q = 3;}
\end{itemize}
the optional input is $tauchenoptions$
\begin{itemize}
 \item Can either not parallelize (\textcolor{RoyalBlue}{tauchenoptions.parallel=0}), or parallelize on graphics card \\ (\textcolor{RoyalBlue}{tauchenoptions.parallel=2}).
    (If you set \textcolor{RoyalBlue}{tauchenoptions.parallel=1} it will just act as if you set a value of zero, as parallelizing on CPU does not appear to help at all)
\end{itemize}
The outputs are,
\begin{itemize}
 \item The discrete grid (of $znum$ points). A column vector. \\
    \textcolor{RoyalBlue}{z\_{}grid}
 \item The transition matrix; the element in row $i$, column $j$ gives the probability of moving to state $j$ tomorrow, given that today is state $i$, \\
    \textcolor{RoyalBlue}{pi\_{}z} \\
    (each row sums to one; size is $znum$-by-$znum$)
\end{itemize}






